\documentclass[12pt]{report}
\usepackage[a4paper,margin=2.5cm]{geometry}
\usepackage{booktabs,tabularx,array,ragged2e}
\begin{document}
\chapter{Results}
\section{Classical modal baseline against the physics-informed network}
Five localisation scores are reported across this chapter, each on a different subset of the
campaign. Table~\ref{tab:conventions} states the basis of each.

\begin{table}[htbp]
\centering
\small
\caption[Localisation scoring conventions]{Localisation scoring conventions. Each analysis
uses the records its inputs permit, and the denominators differ accordingly.}
\label{tab:conventions}
\begin{tabularx}{\textwidth}{@{}
    >{\RaggedRight\arraybackslash\hsize=0.80\hsize}X
    r
    >{\RaggedRight\arraybackslash\hsize=1.20\hsize}X @{}}
\toprule
Analysis & Cases & Basis \\
\midrule
Fixed two-mode signature space $(f_1,f_3)$
  & 11 & Excludes the one Floor 3 replicate with no resolved $\Delta f_3$ \\
Fixed three-mode signature space
  & 9 & Additionally excludes all three Floor 3 replicates, whose $f_2$ is voided by the
      second harmonic of $f_1$ \\
Sensor relocation
  & 18 & Three locations $\times$ three replicates $\times$ two alternative sensor positions,
      scored against top-storey references \\
Classical method against the network
  & 6 & Severe replicates at Floor 1 and Floor 2. The base plate is excluded because $k_1$ is
      the only call the network's output space permits there, and Floor 3 has no severe cell \\
Tap scatter as a location feature
  & 24 & Replicate-level records across all four locations and four grades, reported as an
      exploratory feature outside the main pipeline \\
\bottomrule
\end{tabularx}
\end{table}
\end{document}
